\section{Derivations for the Red Identity Rule}
\label{app:ir-rule-derivation}

Throughout this appendix, we work in
\[
\zxopt=\zxv\setminus\{(I_r)\}
=\{(S),(I_g),(E),(CP),(B),(EU),(HD),(H)\}
+\text{ connectivity}.
\]

The following derivation of \((I_r)\) is highly technical, and most of the difficulty comes from the critical tracking of scalars and the scalar cleanup rule in \hyperref[eq:ir-derivation-L7]{Lemma~\ref*{lem:ir-derivation-L7}}.

We first derive two useful identities from the ruleset:

\newcommand{\irshortstatement}[3]{%
\begin{equation*}
\vcenter{\hbox{\hkstikzitinput{HKS/figures/ir-derivation/renaud-short/#1}}}
\tag{\ensuremath{#3}}
\label{eq:ir-derivation-#2}
\end{equation*}%
}
\newcommand{\irshortrow}[1]{%
\leavevmode\par
\nopagebreak[4]
\begin{center}
\makebox[\linewidth][c]{#1}%
\end{center}
}
\newcommand{\irshortjoin}{\hspace{1em}}
\newcommand{\scalarpairnodes}[3]{%
\begin{scope}[shift={(#2,#3)}]
	\node [style=X dot] (#1-chi-x) at (-0.6,0.35) {};
	\node [style=Z dot] (#1-chi-z) at (-0.6,-0.35) {};
	\node [style=X phase dot] (#1-omega-x) at (0.6,0.45) {$\pi$};
	\node [style=Z dot] (#1-omega-z) at (0.6,-0.45) {};
\end{scope}%
}
\newcommand{\scalarpairedges}[1]{%
\draw (#1-chi-x) to (#1-chi-z);
\draw [bend left=45] (#1-chi-x) to (#1-chi-z);
\draw [bend right=45] (#1-chi-x) to (#1-chi-z);
\draw (#1-omega-x) to (#1-omega-z);
}

\Needspace{0.30\textheight}
\noindent
\begin{minipage}[t]{0.48\linewidth}
\vspace{0pt}
\begin{lemma}
\label{lem:ir-derivation-L1}
\irshortstatement{ir-l01-statement.tikz}{L1}{\mathrm{H}_0}
\end{lemma}
\begin{irproof}
\widehkstikzitinput{HKS/figures/ir-derivation/renaud-short/ir-l01-proof.tikz}
\vspace{-2.6\baselineskip}
\end{irproof}
\end{minipage}\hfill
\begin{minipage}[t]{0.48\linewidth}
\vspace{0pt}
\begin{lemma}
\label{lem:ir-derivation-L2}
\irshortstatement{ir-l02-statement.tikz}{L2}{\mathrm{HC}}
\end{lemma}
\begin{irproof}
\widehkstikzitinput{HKS/figures/ir-derivation/renaud-short/ir-l02-proof.tikz}
\end{irproof}
\end{minipage}

\vspace{\baselineskip}

Now, we derive a rule which lets us remove the red identity along with the two scalar factors shown:

\Needspace{0.35\textheight}
\begin{lemma}
\label{lem:ir-derivation-L3}
\irshortstatement{ir-atomic-l03-statement.tikz}{L3}{\mathrm{LI}_r}
\end{lemma}
\begin{irproof}
\irshortrow{%
	\hkstikzitinput{HKS/figures/ir-derivation/renaud-short/ir-atomic-l03-proof-a.tikz}\irshortjoin
	\hkstikzitinput{HKS/figures/ir-derivation/renaud-short/ir-atomic-l03-proof-b.tikz}\irshortjoin
	\hkstikzitinput{HKS/figures/ir-derivation/renaud-short/ir-atomic-l03-proof-b2.tikz}%
}
\end{irproof}

The Hadamard is self-inverse up to a scalar residue:

\Needspace{0.25\textheight}
\begin{lemma}
\label{lem:ir-derivation-L4}
\irshortstatement{ir-atomic-l04-statement.tikz}{L4}{\mathrm{H}^{-1}}
\end{lemma}
\begin{irproof}
\widehkstikzitinput{HKS/figures/ir-derivation/renaud-short/ir-atomic-l04-proof.tikz}
\end{irproof}

We derive a scalar-tracked version of red spider fusion:

\Needspace{0.35\textheight}
\begin{lemma}
\label{lem:ir-derivation-L5}
\irshortstatement{ir-atomic-l07-statement.tikz}{L5}{\mathrm{LS}_r}
\end{lemma}
\begin{irproof}
\widehkstikzitinput{HKS/figures/ir-derivation/renaud-short/ir-atomic-l07-proof-a.tikz}
\vspace{-0.75\baselineskip}
\irshortrow{%
	\hkstikzitinput{HKS/figures/ir-derivation/renaud-short/ir-atomic-l07-proof-a2.tikz}\irshortjoin
	\hkstikzitinput{HKS/figures/ir-derivation/renaud-short/ir-atomic-l07-proof-b.tikz}%
}
\vspace{-1.5\baselineskip}
\end{irproof}

\Needspace{0.45\textheight}
We obtain a useful scalar-tracked form of the familiar Hopf law:

\begin{lemma}
\label{lem:ir-derivation-L6}
\irshortstatement{ir-renaud-l06-statement.tikz}{L6}{\mathrm{HF}}
\end{lemma}
\begin{irproof}
\irshortrow{%
	\hkstikzitinput{HKS/figures/ir-derivation/renaud-short/ir-renaud-l06-proof-1.tikz}\irshortjoin
	\hkstikzitinput{HKS/figures/ir-derivation/renaud-short/ir-renaud-l06-proof-2.tikz}%
}
\irshortrow{%
	\hkstikzitinput{HKS/figures/ir-derivation/renaud-short/ir-renaud-l06-proof-3.tikz}\irshortjoin
	\hkstikzitinput{HKS/figures/ir-derivation/renaud-short/ir-renaud-l06-proof-a2.tikz}%
}
\irshortrow{%
	\hkstikzitinput{HKS/figures/ir-derivation/renaud-short/ir-renaud-l06-proof-a3.tikz}\irshortjoin
	\hkstikzitinput{HKS/figures/ir-derivation/renaud-short/ir-renaud-l06-proof-6.tikz}\irshortjoin
	\hkstikzitinput{HKS/figures/ir-derivation/renaud-short/ir-renaud-l06-proof-7.tikz}%
}
\end{irproof}

We now obtain the scalar cleanup rule which lets us remove the following pair of scalar factors:

\Needspace{0.45\textheight}
\begin{lemma}
\label{lem:ir-derivation-L7}
\irshortstatement{ir-atomic-l14-statement.tikz}{L7}{\mathrm{IV}'}
\end{lemma}
\begin{irproof}
\irshortrow{%
	\hkstikzitinput{HKS/figures/ir-derivation/renaud-short/ir-renaud-l07-proof-1.tikz}\irshortjoin
	\hkstikzitinput{HKS/figures/ir-derivation/renaud-short/ir-renaud-l07-proof-3.tikz}\irshortjoin
	\hkstikzitinput{HKS/figures/ir-derivation/renaud-short/ir-renaud-l07-proof-4.tikz}%
}
\irshortrow{%
	\hkstikzitinput{HKS/figures/ir-derivation/renaud-short/ir-renaud-l07-proof-5.tikz}\irshortjoin
	\hkstikzitinput{HKS/figures/ir-derivation/renaud-short/ir-renaud-l07-proof-6.tikz}\irshortjoin
	\hkstikzitinput{HKS/figures/ir-derivation/renaud-short/ir-renaud-l07-proof-6a.tikz}%
}
\irshortrow{%
	\hkstikzitinput{HKS/figures/ir-derivation/renaud-short/ir-renaud-l07-proof-7.tikz}\irshortjoin
	\hkstikzitinput{HKS/figures/ir-derivation/renaud-short/ir-renaud-l07-proof-7a.tikz}\irshortjoin
	\hkstikzitinput{HKS/figures/ir-derivation/renaud-short/ir-renaud-l07-proof-8.tikz}\irshortjoin
	\hkstikzitinput{HKS/figures/ir-derivation/renaud-short/ir-renaud-l07-proof-9.tikz}%
}
\end{irproof}

Two necessary technical lemmas:

\Needspace{0.45\textheight}
\begin{lemma}
\label{lem:ir-derivation-L8}
\irshortstatement{ir-renaud-l08-statement.tikz}{L8}{\mathrm{IV}''}
\end{lemma}
\begin{irproof}
\irshortrow{%
	\hkstikzitinput{HKS/figures/ir-derivation/renaud-short/ir-renaud-l08-proof-1.tikz}\irshortjoin
	\hkstikzitinput{HKS/figures/ir-derivation/renaud-short/ir-renaud-l08-proof-2.tikz}\irshortjoin
	\hkstikzitinput{HKS/figures/ir-derivation/renaud-short/ir-renaud-l08-proof-3.tikz}\irshortjoin
	\hkstikzitinput{HKS/figures/ir-derivation/renaud-short/ir-renaud-l08-proof-4.tikz}%
}
\irshortrow{%
	\hkstikzitinput{HKS/figures/ir-derivation/renaud-short/ir-renaud-l08-proof-5.tikz}\irshortjoin
	\hkstikzitinput{HKS/figures/ir-derivation/renaud-short/ir-renaud-l08-proof-6.tikz}\irshortjoin
	\hkstikzitinput{HKS/figures/ir-derivation/renaud-short/ir-renaud-l08-proof-7.tikz}%
}
\irshortrow{%
	\hkstikzitinput{HKS/figures/ir-derivation/renaud-short/ir-renaud-l08-proof-7a.tikz}\irshortjoin
	\hkstikzitinput{HKS/figures/ir-derivation/renaud-short/ir-renaud-l08-proof-8a.tikz}\irshortjoin
	\hkstikzitinput{HKS/figures/ir-derivation/renaud-short/ir-renaud-l08-proof-9.tikz}%
}
\end{irproof}

\Needspace{0.40\textheight}
The following lemma is a useful identity allowing the red state to delete a green phase gate of \(-\pi/2\):

\begin{lemma}
\label{lem:ir-derivation-L9}
\irshortstatement{ir-renaud-l09-statement.tikz}{L9}{\mathrm{D}_{-}}
\end{lemma}
\begin{irproof}
\irshortrow{%
	\hkstikzitinput{HKS/figures/ir-derivation/renaud-short/ir-renaud-l09-proof-a.tikz}\irshortjoin
	\hkstikzitinput{HKS/figures/ir-derivation/renaud-short/ir-renaud-l09-proof-b.tikz}%
}
\irshortrow{%
	\hkstikzitinput{HKS/figures/ir-derivation/renaud-short/ir-renaud-l09-proof-b2.tikz}\irshortjoin
	\hkstikzitinput{HKS/figures/ir-derivation/renaud-short/ir-renaud-l09-proof-b3.tikz}\irshortjoin
	\hkstikzitinput{HKS/figures/ir-derivation/renaud-short/ir-renaud-l09-proof-b4.tikz}\irshortjoin
	\hkstikzitinput{HKS/figures/ir-derivation/renaud-short/ir-renaud-l09-proof-c.tikz}%
}
\end{irproof}

A useful identity is derived:

\Needspace{0.35\textheight}
\begin{lemma}
\label{lem:ir-derivation-L10}
\irshortstatement{ir-renaud-l10-statement.tikz}{L10}{\mathrm{D}'_{\pi}}
\end{lemma}
\begin{irproof}
\irshortrow{%
	\hkstikzitinput{HKS/figures/ir-derivation/renaud-short/ir-renaud-l10-proof-a.tikz}\irshortjoin
	\hkstikzitinput{HKS/figures/ir-derivation/renaud-short/ir-renaud-l10-proof-b.tikz}\irshortjoin
	\hkstikzitinput{HKS/figures/ir-derivation/renaud-short/ir-renaud-l10-proof-c.tikz}%
}
\irshortrow{%
	\hkstikzitinput{HKS/figures/ir-derivation/renaud-short/ir-renaud-l10-proof-d.tikz}\irshortjoin
	\hkstikzitinput{HKS/figures/ir-derivation/renaud-short/ir-renaud-l10-proof-e.tikz}%
}
\end{irproof}

\Needspace{0.32\textheight}
We reproduce the derivation of \((I_r)\) below for readability.

\begin{theorem}
The \((I_r)\) rule is derivable:
\begin{equation*}
\zxopt\vdash
\vcenter{\hbox{\hkstikzitinput{HKS/figures/ir-derivation/ir-atomic-ir-statement.tikz}}}
\end{equation*}
\end{theorem}

\begin{irproof}
\leavevmode\par
\begin{center}
\makebox[\linewidth][c]{%
	\hkstikzitinput{HKS/figures/ir-derivation/renaud-short/ir-renaud-ir-proof-a.tikz}\hspace{1em}%
	\hkstikzitinput{HKS/figures/ir-derivation/renaud-short/ir-renaud-ir-proof-b.tikz}\hspace{1em}%
	\hkstikzitinput{HKS/figures/ir-derivation/renaud-short/ir-renaud-ir-proof-c.tikz}%
}
\end{center}
\end{irproof}
